\section{Discussion}\label{sec:discussion}

\subsection{Accuracy summary}

Table~\ref{tab:results} summarizes the prediction accuracy before
and after applying the corrections described in Section~\ref{sec:main}.

\begin{table}[ht]
\centering
\caption{Prediction accuracy across the four benchmark splits.
Best results in each column are in \textbf{bold}.}\label{tab:results}
\begin{tabular}{@{}lcccc@{}}
\toprule
Split & Baseline & Improved & $\Delta$ & Instances \\
\midrule
\textsc{Normal} & 71.0\% & \textbf{82.4\%} & $+11.4$ & 1000 \\
\textsc{Hard1}  & 65.2\% & \textbf{73.9\%} & $+8.7$  & 69 \\
\textsc{Hard2}  & \textbf{99.5\%} & 95.5\% & $-4.0$  & 200 \\
\textsc{Hard3}  & 52.75\% & \textbf{63.5\%} & $+10.75$ & 400 \\
\midrule
Overall         & 69.8\% & \textbf{79.1\%} & $+9.3$  & 1669 \\
\bottomrule
\end{tabular}
\end{table}

The improvements on \textsc{Normal}, \textsc{Hard1}, and
\textsc{Hard3} are substantial.  The regression on \textsc{Hard2}
is discussed in Section~\ref{subsec:hard2}.

\subsection{Remaining error analysis on \textsc{Hard3}}

After all corrections, $146$ errors remain on the $400$-instance
\textsc{Hard3} split.  These fall into four categories:
\begin{enumerate}[label=(\roman*)]
\item \textbf{Default false negatives (36 instances):} Pairs that
  fall through all rules and receive the default prediction of
  \textsc{false}.  These are predominantly bare source equations
  with $\var(E_1) \in \{2, 3\}$ and non-bare source equations that
  require algebraic reasoning beyond structural features.

\item \textbf{Same-variable-count false negatives (32 instances):}
  Pairs where $\var(E_1) = \var(E_2)$ and a rule correctly predicts
  \textsc{false} in most cases, but the specific pair happens to be
  a true implication.  Resolving these cases requires algebraic
  analysis of the specific equations involved.

\item \textbf{Residual motif false positives (31 instances):}
  Cases where the guarded contradiction motifs $C_{10}$ and $C_{13}$
  still fire incorrectly.  Further tightening of guard conditions
  risks reducing the true-positive count below acceptable levels.

\item \textbf{Size-based false negatives (10 instances):} Cases
  where $\siz(E_2) < \siz(E_1)$ but $E_1 \models E_2$ holds.
  After removing rule~$F_2$, a residual effect comes from other
  rules that incorporate size comparisons.
\end{enumerate}

\begin{remark}\label{rem:algebraic-gap}
Categories (i) and~(ii) together account for $68$ of $146$ errors
($46.6\%$) and represent a fundamental limitation of structural
prediction.  These cases require either direct algebraic proof (via
term rewriting or equational deduction) or counterexample construction
(via finite magma search), neither of which can be reliably encoded
in a small set of syntactic rules.
\end{remark}

\subsection{The \textsc{Hard2} regression}\label{subsec:hard2}

The \textsc{Hard2} split decreased from $99.5\%$ to $95.5\%$,
a loss of $9$ correctly classified pairs.  All $9$ new errors
are false negatives caused by the removal of rules $C_8$, $C_{14}$,
$T_5$, and $T_7$.  These rules happened to fire correctly on
\textsc{Hard2} instances despite being globally unreliable.

This regression is an inherent tension in the multi-split
evaluation framework.  The \textsc{Hard2} split was designed so
that standard structural rules classify most pairs correctly,
while \textsc{Hard3} was designed to defeat exactly those rules.
Removing rules that \textsc{Hard3} exploits necessarily harms
\textsc{Hard2} performance on the marginal cases where those
rules were coincidentally correct.

The trade-off is quantitatively favorable: the $9$ new errors on
\textsc{Hard2} are outweighed by $43$ new correct predictions on
\textsc{Hard3} (from $211$ to $254$ correct out of $400$) and
$114$ new correct predictions on \textsc{Normal}.

\subsection{Log-loss analysis}

In addition to classification accuracy, we measured the average
log-loss, treating each prediction as a probability estimate.
Table~\ref{tab:logloss} reports the results.

\begin{table}[ht]
\centering
\caption{Average log-loss by split (lower is better).
Best results in \textbf{bold}.}\label{tab:logloss}
\begin{tabular}{@{}lcc@{}}
\toprule
Split & Baseline & Improved \\
\midrule
\textsc{Normal} & $-1.343$ & $\mathbf{-0.819}$ \\
\textsc{Hard1}  & $-1.608$ & $\mathbf{-1.209}$ \\
\textsc{Hard2}  & $\mathbf{-0.033}$ & $-0.217$ \\
\textsc{Hard3}  & $-2.181$ & $\mathbf{-1.687}$ \\
\bottomrule
\end{tabular}
\end{table}

The log-loss reductions of $39\%$ (\textsc{Normal}), $25\%$
(\textsc{Hard1}), and $23\%$ (\textsc{Hard3}) indicate that
the improved rules not only classify more instances correctly
but do so with higher confidence on average.

\subsection{Connection to Stone pairings}

Berlioz and Melli\`es~\cite{berlioz2026latent} showed that
implications in the equational theory lattice tend to follow a
``mainstream direction'' in the Stone pairing space: from equations
with lower satisfaction probability to equations with higher
satisfaction probability.  Hard implications---those that are
difficult for automated methods---are precisely the
``contrarian'' edges that go against this flow.

Our structural predictor implicitly captures some aspects of
this geometry.  For instance, bare source equations with many
variables have low satisfaction probability on random magmas
(since all variable assignments must satisfy the identity),
and the corresponding implication targets tend to have higher
satisfaction probability.  The projection collapse lemmas
(Theorems~\ref{thm:left-proj} and~\ref{thm:right-proj})
correspond to equations at extreme points in the Stone pairing
space: the left- and right-projection identities have
satisfaction probability $1/n$ and $1/n$ respectively on a
random $n$-element magma.

However, our predictor cannot capture the full geometry of the
Stone pairing space, which requires computing satisfaction
probabilities on specific finite magmas.  This suggests that a
hybrid approach---combining structural rules with a small table
of Stone pairing values for canonical magmas---could yield
further improvements.

\subsection{Comparison with computational approaches}

The Equational Theories Project~\cite{bolan2025equational} used
automated theorem provers (Vampire~\cite{kovacs2013vampire},
Prover9~\cite{mccune2003prover9}) and systematic counterexample
search to resolve all implications.  A BFS-based rewriting predictor
from a prior iteration achieves $72\%$ accuracy on \textsc{Hard3},
compared to our $63.5\%$.  The $8.5$ percentage-point gap arises
from cases requiring:
\begin{itemize}[leftmargin=*,itemsep=2pt]
\item long rewriting chains ($>10$ steps) that structural features
  cannot predict,
\item counterexample construction using magmas of order $\geq 4$
  that cannot be encoded in a $10$\,KB cheatsheet, and
\item algebraic reasoning about absorbing elements and idempotent
  subalgebras.
\end{itemize}
These limitations are inherent to the structural approach and
motivate the hybrid methods discussed in Section~\ref{sec:conclusion}.
